\scanpage{162}
\mspage{82}

% slanted
\margin{autom strict d'un graphe à homoth.}

\[
G_0\ \text{ou}\ G_{A_0}=\overbrace{\operatorname{Aut}_Z(S,A_0,Z)}^{\text{autom des gr. à homothéties}}=\left\{g\in\unsure{\mathfrak S|\mathscr S}_S\,\middle|\,
\begin{array}{l}g\text{ autom. de }(S,A_0)\ \text{i.e. }g(A_0)=A_0\\g\text{ commute à }Z\end{array}\right\}\tag*{(7)}
\]

% slanted
\margin{autom strict d'un multigraphe à homoth.}

\[
\begin{aligned}
G&=\underbrace{\operatorname{Aut}_Z(S,\pi)}_{\substack{\text{autom des}\\\text{multigraphe à}\\\text{homothéties}}}
=\left\{g\in\mathfrak S_S\,\middle|\,\begin{array}{l}g\text{ autom de }(S,\pi)\ \text{i.e. }g(\pi)=\pi\\g\ \deleted{\unsure{\mathrm{se}}}\text{ commute à }Z\end{array}\right\}\\
&\hspace{1em}\text{i.e. l'opération de $g_\pi$ sur $\pi$ se fait par translation}\\
&\hspace{1em}\text{par un \unsure{él.|élt.} de $Z$}\\
&=\left\{g\in\mathfrak S_S\,\middle|\,\begin{array}{l}g\text{ commute à }Z\\g(A_0)\in\pi\end{array}\right\}
\end{aligned}\tag*{(8)}
\]

\[ G_{A_i}\quad\text{indépendant de } i\in\pi \tag*{(9)} \]

% slanted
\margin{\uline{NB} Il faudrait quand même expliciter ici des cas où la structure \unsure{a|à} $(S,\pi)$ implique déjà la \unsure{connaiss|connaissance} $\unsure{G|\text{(9)}}$, ou m$^{\rm s}$ de $Z$, et \uline{aussi} \deleted{\unsure{suffit}} des cas où tout \unsure{autom} de $(S,\pi)$ appartient déjà à \unsure{$G$|$\operatorname{Aut}G$} i.e induit l'autom. \unsure{trivial|triviale} de $Z$ (c'est le cas \unsure{ds|dans} \unsure{les situations|la situation} \unsure{congruentielles|congruentielle} particulièrement\ldots) : il suffit pour ceci que \uline{tout} autom de $(S,A_0)$ commute à $Z$ i.e soit dans $G_0$ -- ce qui est le cas dans les cas que nous avons à envisager, semble-t-il\ldots\ C'est le cas \unsure{sûrement|évidemment} dans le cas \unsure{prépolyédral régulier|des prépolyèdres réguliers} -- $\unsure{G}$ dont il sera fréquemment \unsure{question}\ldots\ Et $\bullet$ \unsure{peut} décrit $G$ en termes de $(S,F)$, ou $(S,\pi)$, on retrouve $Z$ comme le \uline{centre} de $G$. En termes d'un $(S,F_0)$ ou $(S,A_0)$ déjà donné, \deleted{\unsure{on}} mais sans s'assigner \deleted{\unsure{encore}} $\pi$, le ``meilleur choix'' d'un $Z$ qui se propose (peut-on \unsure{être} tout-à-fait valable ds le cas d'un 2-polyèdre \unsure{str$^{\rm t}$|tr.|str.} \uline{régulier}) c'est de prendre $Z=\operatorname{Comm}_{\mathfrak S_S}(G_0)$\ldots}

diagramme

\diagram{\begin{tikzcd}[ampersand replacement=\&,column sep=small,row sep=small]
(10)\quad 1\arrow[r]\&G_0\arrow[r,hook]\&G\arrow[r]\&\overset{\deleted{\unsure{Z'}}}{Z/\unsure{\mathfrak z|\zeta}}\arrow[r,dash]\&1.\\
\&\&Z\arrow[u,hook]\arrow[ur,swap,"{\lambda\mapsto\lambda^2}"]\&\&
\end{tikzcd}}

\[
\begin{aligned}
\mathfrak z&=\operatorname{Ker}(\deleted{\unsure{T}}\lambda\mapsto T_\lambda:Z\longrightarrow\mathfrak S_\pi)\\
&=\{\lambda\in Z\mid\mathrm{Et}_{A_i}(s)=\mathrm{Et}_{A_i}(\lambda s)\quad\forall s\in S,\}\ \forall i\in\pi\rbrace
\end{aligned}\tag*{(11)}
\]

\uline{NB} $\mathfrak z=\{1\}$ si tout sommet est déterminé par ses étoiles\ldots

NB Soit \deleted{\uline{NB} Soit \unsure{$A$}}

\[ \unsure{\mathcal E|\mathscr E}_{A_0}=\{E\in\unsure{\mathcal P|\mathfrak P}(S)\mid\exists s\in S;\ E=\mathrm{Et}_{A_0}(s)\} \tag*{(12)} \]

alors on a

\[ \mathcal E_{A_i}\ \text{indépendant de }i\in\pi,\quad\text{soit }\mathcal E=\mathcal E_\pi. \tag*{(13)} \]

L' \deleted{étoiles} ensemble $\mathcal E$ \inserted{$\subset\mathcal P(S)$} des étoiles pour les diverses structures de graphe $A_i$, mais le \deleted{leurs} ``centres'' \deleted{varient en \unsure{f$^{\rm t}$ d'un}} (ou les centres) d'une étoile fixée $E\subset S$ varie\unsure{(nt]|(nt)} avec la structure de graphe envisagée. Pour tt $g\in G_0$ on a $g(\mathcal E)=\mathcal E$ i.e $g$ transforme étoile en étoile. C'est le cas en particulier pour les \deleted{$\unsure{g}$} $\lambda\in Z$.

% archive sheet number: 161
